
\documentclass{beamer}
\usetheme{limhp}



\usepackage{bbding}
\usepackage{listings}
\usepackage{textcomp}
%\usepackage{ccicons}
\usepackage{booktabs,colortbl}
\usepackage{color, xcolor}
\usepackage{verbatim}
\usepackage{ragged2e} % for justifying
%\usepackage[timeinterval=60]{tdclock} 

%define different colors
\definecolor{codegreen}{rgb}{0,0.6,0}
\definecolor{codegray}{rgb}{0.5,0.5,0.5}
\definecolor{codepurple}{rgb}{0.58,0,0.82}
\definecolor{backcolour}{rgb}{0.95,0.95,0.92}
\definecolor{lgray}{gray}{0.95}
\definecolor{links}{HTML}{F58607}

%for code colours
\lstdefinestyle{mystyle}{
	backgroundcolor=\color{backcolour},   
	commentstyle=\color{codegreen},
	keywordstyle=\color{magenta},
	numberstyle=\tiny\color{codegray},
	stringstyle=\color{codepurple},
	basicstyle=\ttfamily\footnotesize,
	breakatwhitespace=false,         
	breaklines=true,                 
	captionpos=b,                    
	keepspaces=false,                 
	numbers=left,                    
	numbersep=5pt,                  
	showspaces=false,                
	showstringspaces=false,
	showtabs=false,                  
	tabsize=4
}

\lstset{style=mystyle}


\hypersetup{colorlinks,linkcolor=,urlcolor=links}
\newcommand{\pt}{\textbf{\textsf{Python }}}
\newcommand{\ju}{\textbf{\textsf{JUPYTER }}}



\newenvironment<>{varblock}[2][.9\textwidth]{%
	\setlength{\textwidth}{#1}
	\begin{actionenv}#3%
		\def\insertblocktitle{#2}%
		\par%
		\usebeamertemplate{block begin}}
	{\par%
		\usebeamertemplate{block end}%
	\end{actionenv}}



\makeatletter
\renewcommand\insertshorttitle[1][]{%
	\beamer@setupshort{#1}%
	\let\thanks=\@gobble%
	\hyperlinkpresentationstart{{\beamer@shorttitle}}%
	}
\makeatother


\usepackage[absolute,overlay]{textpos} 
\setlength{\TPHorizModule}{1mm}
\setlength{\TPVertModule}{1mm} 
\newcommand{\MyLogo}{
	\begin{textblock}{15}(107.5,9)
		\includegraphics[width=1.5cm]{HTWlogo}
	\end{textblock}
}

%setting itemize design
\useinnertheme{rectangles}
  \setbeamercolor{item projected}{bg=red!60!black,fg=white}
\setbeamertemplate{sections/subsections in toc}[square]
\setbeamertemplate{items}[circle]


%remove navigation symbols
\beamertemplatenavigationsymbolsempty  
%ToC always on page 2
\def\hyperlinkpresentationstart{\hyperlink{Navigation2}}

\title[HM with Jupyter]{\Large\textbf{Computation in Python/JUPYTER}}
\subtitle{\textbf{\Large\textbf{Introduction}}}
\author[P.K. Yadav, HTW/TUD]{\Large \textbf{\underline{Yadav, Prabhas. K.} }}
\institute[]{\textbf{\textbf{\small \href{mailto:prabhaskumar.yadav@htw-dresden.de}{prabhaskumar.yadav@htw-dresden.de}}}}
\date[13. Oct. 2022]{\begin{textblock}{45}(1.5,72)
		\textbf{13 Oct. 2022}
	\end{textblock}} 

	
%setting the figures folder	
\graphicspath{{./figures/}}	

	
 \begin{document}
 \addtobeamertemplate{block begin}{\setlength\abovedisplayskip{0pt}}	
 	
{

	\renewcommand{\insertnavigation}[1]{}
	\setbeamertemplate{footline}{} 
	\usenavigationsymbolstemplate{}
	\begin{frame}
			\MyLogo
				\titlepage
	\end{frame}
	
	
}

{
	 	%	\renewcommand{\insertnavigation}[1]{}
 	\begin{frame}{\textbf{Today's Contents}}
 			\MyLogo
 		\tableofcontents
 	\end{frame}
 	
}
 	
 
 	\section{Introduction to \pt \& \ju}
 	\begin{frame}[fragile]{\textbf{What is Python? }}
 		\MyLogo
\pt is an \underline{open-source} \underline{interpreted}, \underline{high-level}, general-purpose programming language. This means:\\\vskip0.3cm
    \begin{enumerate}
    	\justifying 
\item \textbf{Interpreted/high-level language}: This makes us avoid the nuances of fundamental coding as done by computer programmers/engineers.\vskip0.2cm
\item \textbf{General purpose programming language and Open-source ecosystem}: This means it is extensible. Already over 200,000 \pt packages are available (check \href{https://pypi.org/}{here}). Also, it means  \textbf{FREE} of cost.\vskip0.2cm
\item \textbf{For Hydraulics Modelling}: We can use \pt packages such as \verb|Numpy|  (for numerical computing), \verb|Scipy| (for scientific computing), \verb|Sympy| (for symbolic computing), \verb|Matplotlib| (for plotting) etc. for our computing and modelling in the course.
    \end{enumerate}

 	\end{frame} 
 	
 \begin{frame} \frametitle{\textbf{A Bit of \pt History}}
 	\begin{itemize}
 		 		\MyLogo
 		\justifying
 		\item Guido van Rossum began developing  \pt in 1980 at Centrum Wiskunde \& Informatica (\href{https://www.cwi.nl/}{CWI}), the Netherlands. Its implementation (\pt v.1) was released in 1994.\vskip0.2cm
 		
 		\item \pt 2.0, released in 2000 became one of the most used general purpose programming language. \pt 2.0 is now being replaced by \pt 3.0 (from 2020). \vskip0.2cm
 		
 		\item \pt 3.0 will be used in our class. It is \textbf{not} 100\% compatible with earlier versions of \pt.\vskip0.2cm
 		
 		\item \pt name comes from the British comedy group \href{https://www.youtube.com/watch?v=vnciwwsvNcc}{Monty Python} (Van Rossum enjoyed their show).
 		 The official Python documentation (\href{https://www.python.org/doc/}{check here}) also contains various references to Monty Python routines.
 		
 	\end{itemize}
 	
 \end{frame}	
 	
 
\begin{frame}
	\frametitle{\textbf{Why Use Python?}}
	 		\MyLogo
	\begin{itemize}
		\justifying
		\item \pt is a \textbf{common tool} among engineers, experts and researchers at universities and industry. 
		\vspace{5.5pt}
		\item \pt is s\textbf{ystem independent}, therefore it is highly portable. Beside, it is a versatile (multi-purpose) language.

		\vspace{5.5pt}
		\item \pt is incredibly flexible and can be adapted to specific local needs using enormous number of \textbf{PACKAGES}. Beside, it can easily interface with other languages
		\vspace{5.5pt}
		\item \pt is under incredibly active development, improving greatly, and supported \textbf{wildly} by both professional and academic developers
	\end{itemize}
	
\end{frame}

\begin{frame}\frametitle{\textbf{Popularity of \pt}}
		 		\MyLogo
	\vspace{-0.5cm}
	
	\begin{columns}
		
		\begin{column}{0.6\linewidth}
			
			\includegraphics[width=1.1\textwidth]{fig_P1_0.png}
			
			
		\end{column}
		
		\begin{column}{0.4\textwidth}
			\begin{enumerate}
				\justifying 
				\item \pt has become a mainstream computing language.
				\item Details on plot are \href{https://insights.stackoverflow.com/}{here.}
				\item This all means - it is good to learn to code in \pt
			\end{enumerate}
		\end{column}
		
	\end{columns}
	
\end{frame}
 	
\section{Very basics of Python Language}
\begin{frame}[fragile]\frametitle{Basic of \pt{} language}
	\justifying 	 		\MyLogo
	Very basics of \pt{}before starting \ju Notebook
	
	\begin{minipage}{0.6\linewidth}
		
		\begin{varblock}[10cm]{}
			
			\begin{itemize}
				\item[\HandRight] \textcolor{red}{\textbf{Comments symbol}}: with  \textbf{\#}
				\item[\HandRight] \textcolor{red}{\textbf{End statement}}: with semicolon (;) or line-end 
				\item[\HandRight] \textcolor{red}{\textbf{Multi-line statement}}: with continuation character (\textbf{\textbackslash})
					\item[\HandRight] \textcolor{red}{\textbf{Multiple statement}}: separate each with \textbf{;}
				\item[\HandRight] \textcolor{red}{\textbf{Code Indentation}}: is required (4 spaces for each)
				\item[\HandRight] \textcolor{red}{\textbf{String quotes}}: use either \textbf{`` "} or \textbf{` '}
				\item[\HandRight] \textcolor{red}{\textbf{Naming}}: is case sensitive $b\neq B$
				\item[\HandRight] \textcolor{red}{\textbf{Variable Assignment}}: use \textbf{=}, e.g. \verb|x=5|
				\item[\HandRight] \textcolor{red}{\textbf{Multiple Assignment}}: allowed e.g, \verb|a = b = c = 1| \\ \hskip4cm or \verb|a,b,c = 1,7,10|
				\item[\HandRight] \textcolor{red}{\textbf{Help}}: use \verb|help()| and \textcolor{red}{\textbf{get the cheatsheet from}} \href{https://github.com/prabhasyadav/HTW/blob/master/Python3-cheatsheet.pdf}{\textbf{here}}
			\end{itemize}
			
		\end{varblock}	
	\end{minipage}
	
\end{frame}

\begin{frame}[fragile]\frametitle{Basic of \pt{} language}
	\justifying 	 		\MyLogo
	Data types in \pt{}
	
	\begin{minipage}{0.6\linewidth}
		
		\begin{varblock}[10.2cm]{}
			
			\begin{itemize}
				\item[\HandRight] \textcolor{red}{\textbf{Numbers}}: numbers of any types: \verb|int (e.g., 10),|\\\hskip1cm\verb|long (e.g., 5192L ), float (e.g., 13.4),| \\\hskip1cm\verb|complex (e.g., 2.13j)|
				\item[\HandRight] \textcolor{red}{\textbf{String}}: with quotation mark, e.g., \verb|"john"|
				\item[\HandRight] \textcolor{red}{\textbf{List}}:  are \underline{modifiable} placed in \textbf{[\,]} separated by comma (\textbf{,}). \\\hskip1cm e.g.,  \verb|[1, 3, "apple", 2, "jo"]|
				\item[\HandRight] \textcolor{red}{\textbf{Tuple}}: are \underline{non- modifiable} placed in \textbf{(\,)} separated by \\\hskip1.2cm comma (\textbf{,}).  e.g.,  \verb|(11, 3, "ape", 2, "job")|
				\item[\HandRight] \textcolor{red}{\textbf{Dictionary}}: For table-like data. Placed in \textbf{\{\}} separated by comma (\textbf{,}), e.g.,\\ 
				 \verb|{'name': `john', `code':6734, `dept': `sales'}|
			\end{itemize}
		\end{varblock}	
	\end{minipage}
	
\end{frame}


\begin{frame}[fragile]\frametitle{Basic of \pt{} language}
	\MyLogo
	\justifying
Basic \textbf{operators} in \pt


\begin{minipage}{0.6\linewidth}
	
\begin{varblock}[5.6cm]{Arithmetic Operations}
\begin{tabular}{@{}ccc@{}}
	\toprule
	\textbf{Sym.} & \textbf{Operation} & \textbf{e.g.,} \\ \midrule
	+/- & Add/Subs. & a+b \\
	* & Multiply & a*b \\ 
	/ & Divide & a/b \\
	** & exponent & a**2 = $a^2$ \\ 
	\% & Modulus & Reminder a/b \\ \bottomrule
\end{tabular}
\end{varblock}	
\end{minipage}
\hfill
\begin{minipage}{0.7\linewidth}
	
	\begin{varblock}[4.75cm]{Bitwise/Comparison }
		
\begin{tabular}{@{}ccc@{}}
		\toprule
	\textbf{sym.} & \textbf{Operation} & \textbf{e.g.,} \\ \midrule
	\& & AND & a \& b \\
	$\vert$ & OR & a$\vert$b \\
	== & equal to & a==b \\
	!= & not equal to & a!=b \\
	\textgreater{}/\textless{} & greater/less & a\textgreater{}b 
\end{tabular}
		
	\end{varblock}	
\end{minipage}

\vskip0.3cm
Refer to \pt documentation for complete description. \pt documentation is very extensive and can be obtained from \href{https://www.python.org/doc/}{\textbf{here}}

\end{frame}


\begin{frame}[fragile]\frametitle{Basic of \pt{} language}
	\MyLogo
	\justifying
	A \verb|FUNCTION| in \pt\\\vskip0.1cm
	
	A \textbf{function} in a programming provide an ability to develop a reusable code-block with an option of several operations. So, a \textbf{function} have a (or a set) of \textit{input} and provide an (or a set)  \textit{output}. 
	
	\begin{minipage}{0.6\linewidth}
\begin{varblock}[4cm]{\pt Function e.g. }
	
\begin{lstlisting}[language=Python]
def func1(a,b):

	c = a+b
	d = c*a
	
	return c, d

\end{lstlisting}
		\end{varblock}	
	\end{minipage}
\hfill
	\begin{minipage}{0.4\linewidth}
	\begin{varblock}[6.3cm]{What is it?}
	\textcolor{magenta}{\textbf{def}} in line-1 begins a \verb|Function| block \\
	\textbf{func1} in line 1 is user-defined function name\\
	\textbf{a,b} in line 1 are function \verb|input|\\
	\textbf{c,d} in lines 3-4 is function \verb|operation|\\
	\textcolor{magenta}{\textbf{return}} line-6 is function \verb|output| block 	
	\end{varblock}	
\end{minipage}
\vskip0.2cm
\textbf{Semicolon} (:) in line 1 and \textbf{Indentation} after line 1 are required.\\
\textcolor{magenta}{\textbf{def, return}} are \pt keywords. There are quite few of them.

\end{frame}




\section{The JUPYTER Notebook}

\begin{frame}\frametitle{\textbf{The \ju Notebook}}
	\MyLogo
	\justifying
	We avoid hard-coding and long codes in this course. We use something called \textbf{``JUPYTER notebook''}. \textcolor{red}{\textbf{Get \ju cheatsheet from}} \href{https://github.com/prabhasyadav/HTW/blob/master/Jupyter_Cheatsheet.pdf}{\textbf{here}}
	\vspace{0.4cm}
	
	\centering
	\includegraphics[width=0.8\textwidth]{fig_P1_1}
	
\end{frame}

\begin{frame}\frametitle{\textbf{More on \ju{} notebook}}
	\justifying 	\MyLogo
	\ju{} -- a much better way to work in \pt that allows on a single display: \vskip0.3cm
	
	\begin{itemize}
		\justifying
		\item interface for Python and also many other programming language (R, Java, Julia etc.)
		\item writing and editing of codes (\textit{\underline{we will learn this}})
		\item operation of \pt and its packages (e.g., Numpy)
		\item interactively viewing the code output (i.e., results and plots) 
		\item getting \textbf{help}
	\end{itemize}
	\vskip0.3cm
	The \ju{}notebook runs in the browser (Chrome, Firefox etc.). So, nothing really required to be installed. 
	
	\vskip0.3cm
	
	For better computing and learning, you are encouraged to install \ju system following the instructions from \href{https://jupyter.readthedocs.io/en/latest/install.html}{\textbf{here}}.
	
\end{frame}

\begin{frame}[fragile]\frametitle{\textbf{Interacting with \ju{} notebook}}
	\justifying 	\MyLogo
	\ju notebook runs in a browser. \\\vskip0.2cm
	The first webpage with \ju{}will have http address, \\
	e.g., \verb|http://localhost:8888/tree|. Internet is \textbf{not} required. \\\vskip0.3cm
	\ju{}notebook saves file as \verb|IPYNB| file.\\\vskip0.2cm
	
\begin{center}
		\includegraphics[width=0.65\textwidth]{fig_P1_2}
\end{center}

	The first-page of \ju provides a basic \verb|File-Explorer| of your system - from here it is possible to load available \verb|IPYNB| file or create a new one.

	
\end{frame}




\begin{frame}[fragile]\frametitle{\textbf{Interacting with \ju{} notebook}}
	\justifying 	\MyLogo
\ju notebook has a block-based interface called \textbf{CELL}. \\\vskip0.1cm

Codes are put in the CELL and executed with \textbf{RUN} or using Keyboard buttons \textbf{SHIFT+ENTER}\vskip0.1cm

The first \verb|CELL| is usually for importing \verb|PACKAGES| - such as \verb|Numpy|, \verb|Scipy|, \verb|Matplotlib|

\begin{center}
	\includegraphics[width=0.55\textwidth]{fig_P1_3}
\end{center}

Texts in \textcolor{green}{green -- (imports, from, as)} are \pt keywords. \vskip0.1cm
In this block \verb|numpy| and \verb|matplotlib.pyplot| is imported to \ju  \vskip0.1cm
Also a single function of a package can be imported -- in second line the function \verb|erf| is imported

\end{frame}

\begin{frame}[fragile]\frametitle{\textbf{Interacting with \ju{} notebook}}
	\justifying 	\MyLogo
\vspace{-0.3cm}
	For a \textbf{basic} operation (e.g., as a \verb|CALCULATOR|) in \ju{}notebook:
	\begin{enumerate}
		\item Variables (e.g., \verb|a, b, c|) are first defined.\vskip0.05cm
		\item Operation using \verb|Operators| and variables are stated.\vskip0.05cm
	\item Result is printed when \verb|OUTPUT| variable is put (see last line)
		\item Code is \textbf{executed} with \verb|Run| or using keybboard.
	\end{enumerate}

	\begin{center}
		\includegraphics[width=0.65\textwidth]{fig_P1_4}
	\end{center}
	\justifying 
\textbf{NOTE: }	The interactivity between \textcolor{blue}{In [3]} and \textcolor{red}{Out [3]} in the code-block can be repeated only for the particular cell.
\end{frame}

\begin{frame}[fragile]\frametitle{\textbf{Interacting with \ju{} notebook}}
	\justifying 	\MyLogo
	
	For a complex operation (e.g., defining own \verb|FUNCTION|) in \ju{}notebook:\vskip0.1cm
	We will learn these by solving our own problems- as we progress with the course.\vskip0.1cm
	Here is a basic example how it all works. \vskip0.1cm
	
 \begin{columns}
 	\begin{column}{0.5\textwidth}
 		\textbf{Definition parts}
 		\includegraphics[width=1.1\textwidth]{fig_P1_5}
 	\end{column}
 	\begin{column}{0.5\textwidth} 
 		\textbf{Execution and Result parts}
 		\begin{center}
 			\includegraphics[width=0.8\textwidth]{fig_P1_6}
 		\end{center}
 	\end{column}
 \end{columns}
	


	
\end{frame}

\begin{frame}[fragile]\frametitle{\textbf{Next steps with \pt{}and \ju{}}}
	\justifying 	\MyLogo

\pt is \textbf{HUGE}. Becoming a \pt expert or a programmer is not our \textbf{GOAL}. Here we want to make our \textbf{HYDRAULICS MODELLING} computation easy.	For that a very small part of a much smaller part of \pt called \href{https://scientific-python.org/}{\textbf{SCIENTIFIC PYTHON}} is sufficient.\\\vskip0.1cm

So we focus on following package (mostly \textbf{SELF STUDY})
\\\vskip0.1cm
\begin{enumerate}
\item \verb|Numpy|: For matrix and numerical works (see \href{https://numpy.org/doc/}{Numpy docs})
\item \verb|Matplotlib|: For plots and visualizations (see \href{https://matplotlib.org/users/index.html}{Mat docs})
\item \verb|Scipy|: For computations and modelling (see \href{https://docs.scipy.org/doc/}{Scipy docs})
\item \verb|Sympy|: For symbolic calculations (see \href{https://docs.sympy.org/latest/index.html}{Sympy docs})
\item \verb|Pandas|: For database manipulation (see \href{https://pandas.pydata.org/pandas-docs/stable/}{Pandas docs})
\item \verb|Statsmodels|: For statistical modelling (see \href{	https://www.statsmodels.org/stable/examples/index.html}{docs here})
\item and more as required \ldots
\end{enumerate}

\end{frame}



{
	%\setbeamertemplate{navigation symbols}{}
		\renewcommand{\insertnavigation}[1]{}	
	\setbeamertemplate{footline}{} 
	\begin{frame}[fragile]{\huge \textbf{Python is easy$\ldots$,}}
		\beamertemplatenavigationsymbolsempty
		\Large
		
		\textbf{Let us learn to work in \ju environment}  \\\vspace{0.4cm} 

We will focus on \verb|Numpy|, \verb|Scipy| and \verb|Matplotlib| for our computation and modelling requirements.


\renewcommand{\insertnavigation}[1]{}	
		
		
	\end{frame}
	
}		














 	
 	
 	
 \end{document} 